\section{Introduction}\label{section:Introduction}

Quantum computing offers a route to simulating quantum many-body systems in a way that is difficult to reproduce with classical resources, an idea often traced to early arguments by \textcite{feynman1982simulating} about using controllable quantum systems to emulate nature. In parallel, the modern field of quantum machine learning (QML) studies how parameterized quantum circuits can be used for inference, optimization, and simulation tasks, often within hybrid quantum--classical workflows suitable for noisy intermediate-scale quantum (NISQ) devices \parencite{schuld2018supervised,schuld2021mlqc,conti2024qml,carleo2017nqs,carleo2019mlphys}. 

A central example of such hybrid methods is the Variational Quantum Eigensolver (VQE), which approximates the ground state of a quantum Hamiltonian using a parameterized quantum circuit evaluated on a quantum processor and optimized by a classical algorithm \parencite{peruzzo2014vqe,mcclean2016theory,tilly2022vqe_review}. Variational circuit methods are closely related to other NISQ-era algorithms such as the Quantum Approximate Optimization Algorithm (QAOA) \parencite{farhi2014qaoa,zhou2019qaoa}.

The aim of this project is to study a simplified interacting many-body Hamiltonian, known as the Lipkin model, and to compare classical diagonalization with VQE-based energy estimation across a sequence of increasingly structured test problems. I begin by constructing the computational basis for one and two qubits and studying the action of elementary operators and gates, including the Pauli matrices, Hadamard gate, phase gate, and the controlled-NOT (CNOT) gate. This provides a controlled setting for generating entangled states such as the Bell states and for quantifying bipartite entanglement through reduced density matrices and the von Neumann entropy.

I then analyze a simple two-level Hamiltonian written both as a matrix representation and in terms of Pauli operators, and I study how its eigenvalues and eigenvectors evolve as an interaction parameter is varied. This serves as a baseline for implementing the VQE algorithm and comparing its results with exact classical solutions.

Finally, I extend the analysis to interacting two-qubit systems and to the Lipkin Hamiltonian, originally introduced as a solvable model for testing many-body approximation methods \parencite{lipkin1965}. The Hamiltonian is rewritten in terms of Pauli operators so that expectation values can be evaluated using quantum circuits. Classical diagonalization results are then compared with VQE simulations, allowing us to study how interaction strength influences spectral properties, state mixing, and entanglement.

The remainder of this report is structured as follows. \Cref{section:methods} introduces the quantum formalism, the VQE algorithm, and the Lipkin model. \Cref{sec:results} presents results for each system in order of increasing complexity. \Cref{sec:numerical-stability} discusses numerical precision and the limitations of the variational approach. \Cref{sec:conclusion} summarizes the main findings and outlines directions for future work. Supplementary figures, circuit diagrams, and the Python environment are collected in the Appendix.